\documentclass[conference]{IEEEtran}
\IEEEoverridecommandlockouts

\usepackage{cite}
\usepackage{amsmath,amssymb}
\usepackage{graphicx}
\usepackage{booktabs}
\usepackage{url}
\usepackage[hidelinks]{hyperref}

% ---------------------------------------------------------------------------
% Every result number in this paper comes from numbers.tex, a generated table
% or a generated figure (python -m crisisbench.paper_build). Nothing is typed.
% ---------------------------------------------------------------------------
\IfFileExists{numbers.tex}{% Generated by crisisbench.paper_build. Do not edit.
\newcommand{\boostCgnsde}{3.88}
\newcommand{\boostCgnsde2}{29.15}
\newcommand{\boostDccGarch}{8.78}
\newcommand{\boostGarchFhs}{16.94}
\newcommand{\boostHmmGaussian}{4.79}
\newcommand{\boostIidGaussian}{2.77}
\newcommand{\boostQuantgan}{3.26}
\newcommand{\boostRegimeGaussian}{10.35}
\newcommand{\boostStationaryBootstrap}{21.37}
\newcommand{\boostTPpca}{2.24}
\newcommand{\boostTimegan}{n/p}
\newcommand{\boostTimevae}{n/p}
\newcommand{\corrCgnsde}{2.15}
\newcommand{\corrCgnsde2}{1.04}
\newcommand{\corrDccGarch}{2.01}
\newcommand{\corrGarchFhs}{2.09}
\newcommand{\corrHmmGaussian}{0.16}
\newcommand{\corrIidGaussian}{0.14}
\newcommand{\corrQuantgan}{5.00}
\newcommand{\corrRegimeGaussian}{0.15}
\newcommand{\corrStationaryBootstrap}{0.20}
\newcommand{\corrTPpca}{1.62}
\newcommand{\corrTimegan}{33.22}
\newcommand{\corrTimevae}{12.50}
\newcommand{\nGenerators}{12}
\newcommand{\realKrBoost}{22.49}
\newcommand{\realKrBoostHi}{45.92}
\newcommand{\realKrBoostLo}{-0.55}
\newcommand{\realKrCrisis}{513}
\newcommand{\realKrWindows}{3419}
\newcommand{\realNseBoost}{29.33}
\newcommand{\realNseBoostHi}{56.59}
\newcommand{\realNseBoostLo}{7.92}
\newcommand{\realNseCrisis}{514}
\newcommand{\realNseWindows}{3424}
\newcommand{\realUsBoost}{34.46}
\newcommand{\realUsBoostHi}{60.17}
\newcommand{\realUsBoostLo}{8.85}
\newcommand{\realUsCrisis}{525}
\newcommand{\realUsWindows}{3497}
\newcommand{\scoreBCgnsde}{0.79}
\newcommand{\scoreBCgnsde2}{0.31}
\newcommand{\scoreBDccGarch}{0.79}
\newcommand{\scoreBGarchFhs}{0.82}
\newcommand{\scoreBHmmGaussian}{0.03}
\newcommand{\scoreBIidGaussian}{0.04}
\newcommand{\scoreBQuantgan}{1.38}
\newcommand{\scoreBRegimeGaussian}{0.04}
\newcommand{\scoreBStationaryBootstrap}{0.06}
\newcommand{\scoreBTPpca}{0.44}
\newcommand{\scoreBTimegan}{10.00}
\newcommand{\scoreBTimevae}{4.50}
\newcommand{\scoreCCgnsde}{1.22}
\newcommand{\scoreCCgnsde2}{0.15}
\newcommand{\scoreCDccGarch}{0.62}
\newcommand{\scoreCGarchFhs}{0.54}
\newcommand{\scoreCHmmGaussian}{0.70}
\newcommand{\scoreCIidGaussian}{1.03}
\newcommand{\scoreCQuantgan}{1.09}
\newcommand{\scoreCRegimeGaussian}{0.54}
\newcommand{\scoreCStationaryBootstrap}{0.18}
\newcommand{\scoreCTPpca}{0.90}
\newcommand{\scoreCTimegan}{8.79}
\newcommand{\scoreCTimevae}{7.56}
\newcommand{\tauBCkrFull}{0.29}
\newcommand{\tauBCkrHeld}{-0.47}
\newcommand{\tauBCnseFull}{0.45}
\newcommand{\tauBCnseHeld}{0.35}
\newcommand{\tauBCusFull}{0.60}
\newcommand{\tauBCusHeld}{0.05}
\newcommand{\tauRiskA}{1.00}
\newcommand{\tauRiskB}{-0.20}
\newcommand{\tauRiskC}{0.80}
\newcommand{\tauRiskD}{0.60}
\newcommand{\verdictG1}{FAIL}
\newcommand{\verdictG1v2,posthoc}{PASS}
\newcommand{\verdictG2}{FAIL}
\newcommand{\verdictG2v2,posthoc}{FAIL}
\newcommand{\verdictG3}{PASS}
\newcommand{\verdictG4}{PASS}
\newcommand{\verdictH1}{not supported}
\newcommand{\verdictH2}{supported}
\newcommand{\verdictH2v2,posthoc}{not supported}
\newcommand{\verdictH3}{supported}
\newcommand{\verdictH4}{supported}
\newcommand{\verdictH5}{supported}
\newcommand{\verdictKILL}{not triggered}
}{}
\newcommand{\pendingnum}{\textit{n/a}}
% Fallbacks for macros that exist only after the remaining runs have finished.
\providecommand{\realNseBoost}{\pendingnum}
\providecommand{\realNseBoostLo}{\pendingnum}
\providecommand{\realNseBoostHi}{\pendingnum}
\providecommand{\realNseWindows}{\pendingnum}
\providecommand{\realNseCrisis}{\pendingnum}
\providecommand{\realUsBoost}{\pendingnum}
\providecommand{\realUsBoostLo}{\pendingnum}
\providecommand{\realUsBoostHi}{\pendingnum}
\providecommand{\realKrBoost}{\pendingnum}
\providecommand{\realKrBoostLo}{\pendingnum}
\providecommand{\realKrBoostHi}{\pendingnum}
\providecommand{\nGenerators}{\pendingnum}
\providecommand{\tauBCnseFull}{\pendingnum}
\providecommand{\tauBCnseHeld}{\pendingnum}
\providecommand{\tauBCusFull}{\pendingnum}
\providecommand{\tauBCkrFull}{\pendingnum}
\providecommand{\tauRiskB}{\pendingnum}
\providecommand{\tauRiskC}{\pendingnum}
\providecommand{\boostIidGaussian}{\pendingnum}
\providecommand{\boostTPpca}{\pendingnum}
\providecommand{\boostCgnsde}{\pendingnum}
\providecommand{\vOneBestEpoch}{\pendingnum}
\providecommand{\vOneEpochs}{\pendingnum}

% Missing tables and figures must not break compilation.
\newcommand{\maybetable}[1]{\IfFileExists{#1}{\input{#1}}{\textit{(pending)}}}
\newcommand{\maybefig}[2]{\IfFileExists{#2}{\includegraphics[width=#1]{#2}}{\fbox{\textit{(pending)}}}}

\begin{document}

\title{Correlation Is Not Crisis: Evaluating Synthetic Multi-Asset Returns Under Regime Stress}

\author{\IEEEauthorblockN{Yashovardhan Thopte}
\IEEEauthorblockA{\textit{Department of Computer Engineering} \\
\textit{Fr. Conceicao Rodrigues College of Engineering}\\
Mumbai, India}
\and
\IEEEauthorblockN{Alok Shukla}
\IEEEauthorblockA{\textit{Department of Computer Engineering} \\
\textit{Fr. Conceicao Rodrigues College of Engineering}\\
Mumbai, India}
\and
\IEEEauthorblockN{Aaryan Singh}
\IEEEauthorblockA{\textit{Department of Computer Engineering} \\
\textit{Fr. Conceicao Rodrigues College of Engineering}\\
Mumbai, India}
% \and
% \IEEEauthorblockN{(Project guide: name to be added)}
% \IEEEauthorblockA{\textit{Department of Computer Engineering} \\
% \textit{Fr. Conceicao Rodrigues College of Engineering}\\
% Mumbai, India}
}

\maketitle

\begin{abstract}
Synthetic multi-asset return generators are usually judged by how well they reproduce the unconditional
correlation matrix and a list of stylised facts. Stress testing needs something else: the way cross-asset
dependence changes when markets are under stress. We test, with hypotheses and pass/fail gates frozen before
any generator was scored, whether these two views agree. CrisisBench scores generators on four metric families
(marginal and temporal, static dependence, crisis dependence, distinguishability) on a panel of large Indian
equities, with a crisis label computed by one frozen rule from each window's own returns, so that a generator
is never handed a label. Every error is expressed in units of real sampling noise, obtained from a block
bootstrap of the real data, and the benchmark statistic is checked on simulated panels whose crisis correlation
boost is known. On the real panel the crisis correlation boost is \realNseBoost\% with a 95\% interval of
[\realNseBoostLo, \realNseBoostHi], which is wide.
% RESULT-DEPENDENT: one sentence with the Kendall tau between the family-B and family-C rankings on NSE
% (\tauBCnseFull, \tauBCnseHeld) and the H1 verdict, once all generators (neural baselines, CG-NSDE) are scored.
The rank agreement between correlation fidelity and crisis dependence across \nGenerators\ generators is
reported with its uncertainty, on India and on two replication markets.
% RESULT-DEPENDENT: one sentence on whether the family-C ranking agrees with the out-of-sample VaR/ES backtest
% ranking better than the family-B ranking does (H3, gate G3).
We validate the rankings against an out-of-sample Value-at-Risk and Expected-Shortfall backtest.
% RESULT-DEPENDENT: one sentence on where CG-NSDE ranks and whether a regime-conditional Gaussian matches it (H2, G1, G2).
We also describe CG-NSDE, a graph-conditioned latent neural SDE used as a reference method, and report its
ablations and every failed gate. The controlled check shows that the boost statistic is compressed and is not
zero on a GARCH process without regimes, so it must be read as an ordinal quantity.
\end{abstract}

\begin{IEEEkeywords}
synthetic financial data, generative models, benchmark, crisis correlation, stress testing, Value-at-Risk,
Expected Shortfall, neural SDE
\end{IEEEkeywords}

% ===========================================================================
\section{Introduction}
% ===========================================================================
Generative models of asset returns are proposed for data augmentation, privacy-preserving data sharing and
scenario generation for risk management \cite{assefa2020synthetic,potluru2023synthetic,eckerli2021gans}. In the
multi-asset case the usual evidence of quality is a small distance between the generated and the real
correlation matrix, together with stylised facts such as heavy tails and volatility clustering
\cite{cont2001stylized}. A risk manager who wants stressed scenarios asks a different question. Correlations
between equities rise in falling and turbulent markets \cite{longin2001extreme,ang2002asymmetric}, and it is
this change of dependence, not its long-run average, that removes diversification when it is needed. A
generator can match the pooled correlation matrix closely and still produce the same dependence in every
window: an i.i.d. Gaussian with the sample covariance does exactly that.

This paper asks whether the two views give the same ordering of generators. The question is empirical, and we
treat it as a hypothesis to be tested instead of a premise. Before any generator was scored we froze, in one
protocol file, the windows, the splits, the crisis labelling rule, the metric families, the uncertainty
procedure, five hypotheses (H1 to H5) and four gates (G1 to G4) with a rule for abandoning the benchmark
framing. All verdicts in this paper are produced by code from the run outputs, and the failed ones are
reported next to the passed ones.

The design responds to a flaw in our own earlier evaluation. In the previous version of this project each
generated window was conditioned on a real window and inherited that window's crisis label. Crisis metrics
computed in this way partly measure the conditioning data and not the generator, and the earlier headline
numbers could not be reproduced. Here sampling is unconditional for every generator and the crisis label of a
generated window is computed from the generated returns by the same frozen rule that labels real windows.

Our contributions, in order of importance, are as follows.
\begin{enumerate}
\item \textbf{An empirical test.} We test whether correlation-fidelity metrics and crisis-dependence metrics
rank generators differently (H1), whether the difference replicates on other markets (H4), and whether either
ranking agrees with an out-of-sample VaR/ES backtest (H3).
\item \textbf{CrisisBench.} A regime-labelled benchmark in which every error is normalised by a
real-versus-real noise ceiling, with a held-out crisis split, a controlled-ground-truth check of the main
crisis statistic, a data-contamination check (H5) and a downstream risk validation.
\item \textbf{CG-NSDE as a reference method.} A graph-conditioned latent neural SDE with a low-rank Student-t
emission, sampled unconditionally, with component ablations. It is one entry in the benchmark and not the
subject of the paper.
\end{enumerate}

To our knowledge, no existing benchmark of synthetic return generators evaluates cross-asset dependence
separately in crisis and calm regimes, normalises generator error by a real-versus-real noise ceiling, and
checks the resulting ranking against an out-of-sample VaR/ES backtest. Existing evaluations address these
elements separately: regime-wise evaluation \cite{ang2025ctbench}, tail-risk scoring through VaR and ES
\cite{cont2022tailgan,ericson2024var}, and cross-sectional tail co-movement \cite{huh2026marketgans}. This
statement rests on a search of abstracts and not on a systematic review.

% ===========================================================================
\section{Related Work}
% ===========================================================================
\textbf{Generators.} TimeGAN \cite{yoon2019timegan}, Quant GANs \cite{wiese2020quantgans} and TimeVAE
\cite{desai2021timevae} are widely used neural baselines for time-series generation; signature-based
conditional GANs \cite{liao2024sigcwgan}, correlation-matrix GANs \cite{marti2020corrgan}, market generators
for small data \cite{buehler2020market} and neural SDEs \cite{li2020scalable,kidger2021neuralsde} extend the
family. Recent multi-asset work targets inter-asset correlation directly
\cite{masi2023correlated,masi2026gandiffusion,huh2026marketgans}. Classical alternatives are GARCH
\cite{bollerslev1986garch}, dynamic conditional correlation \cite{engle2002dcc}, filtered historical simulation
\cite{baroneadesi1999fhs}, the stationary bootstrap \cite{politis1994stationary} and Markov regime switching
\cite{hamilton1989regime}.

\textbf{Evaluation.} General time-series generation benchmarks and surveys of evaluation measures exist
\cite{ang2023tsgbench,stenger2024evaluation}. In finance, Kwon and Lee test whether GANs learn stylised facts
\cite{kwon2024stylized}; Caulfield and Gleeson compare deep generative and parametric models on multivariate
series \cite{caulfield2024systematic}; Cetingoz and Lehalle argue that standard generator evaluation misses
what matters for portfolio use \cite{cetingoz2025synthetic}; Zhang et al. show that visually realistic output
can fail under a trading backtest \cite{zhang2026sfag}; Hounwanou et al. compare ARIMA-GARCH, a VAE and TimeGAN
on a single index \cite{hounwanou2025synthetic}. CTBench evaluates crypto generators over several market
regimes and includes a risk dimension \cite{ang2025ctbench}. Tail-GAN uses the joint elicitability of VaR and ES
as a training objective \cite{cont2022tailgan}; Ericson et al. use VaR backtesting as one performance indicator
on yield-curve data \cite{ericson2024var}; Flaig and Junike build a GAN-based scenario generator as a market
risk model \cite{flaig2022scenario}; MARCD conditions a diffusion model on regimes for allocation
\cite{alzahrani2025marcd}. Our reading of these works is based on their abstracts. None of the abstracts
mentions a real-versus-real noise ceiling or a crisis-labelled held-out test of cross-asset dependence; the
closest in spirit to our thesis is \cite{cetingoz2025synthetic}.

\textbf{Crisis dependence.} Higher correlation in bear markets and in the tails is documented in
\cite{longin2001extreme,ang2002asymmetric}; systemic measures based on the share of variance absorbed by a few
factors and on network connectedness follow the same idea \cite{kritzman2011absorption,diebold2014network}.
Forbes and Rigobon \cite{forbes2002contagion} show that measured correlation rises with volatility even when
the underlying dependence is unchanged. This caveat applies directly to our main crisis statistic and is
examined in Section~\ref{sec:dgp}.

% ===========================================================================
\section{CrisisBench: Benchmark Design}
\label{sec:bench}
% ===========================================================================
\subsection{Data}
The main panel holds daily log returns of 47 large Indian equities from the NIFTY 50 universe, from November
2010 to December 2024, downloaded from Yahoo Finance (auto-adjusted close). Two replication panels of the same
size and period are built by the same code from large-cap United States and Korean (KOSPI) stocks. The cleaning
rules are identical for all markets and were fixed before scoring: a ticker is kept if its history starts at
the beginning of the sample and runs to its end, it has fewer than 2\% missing days and no price run unchanged
for more than ten trading days; gaps of at most five days are forward filled; a return above 30\% in absolute
value is replaced by the cross-sectional median when it reverses by at least 80\% within two days (bad tick) or
exceeds 50\% without reversal (unadjusted corporate action), and is otherwise kept and logged as a real move.
Every action is written to an audit file with a hash of the return matrix.

Two properties of this panel limit what can be concluded. First, the candidate list is a recent constituent
list and a full history is required, so the panel has survivorship bias: firms that left the index or listed
late are absent. Second, Yahoo data is a free vendor feed without point-in-time guarantees; the raw prices
cannot be redistributed, so the panel is reproducible only by re-download.

Stocks on which derivatives trade, which covers the NIFTY 50 constituents, have no fixed daily price band in
India. They have a dynamic band of 10\% that is relaxed in steps of 5\% after a cooling-off period of 15
minutes \cite{sebi2023tradingchapter}. Daily returns of these names are therefore not truncated at a fixed
limit, and we apply no circuit-filter clipping to real or generated returns. Index-level truncation can come
only from the market-wide circuit breaker.

\subsection{Windows, fits and seeds}
A sample is a window of $T=60$ trading days and $N$ assets. Real windows overlap with a stride of one day; the
main panel gives \realNseWindows\ windows, of which \realNseCrisis\ are labelled crisis. Each generator is run
in two fits. The \emph{full} fit estimates on the whole sample and is scored against the whole sample
(in-sample fidelity). The \emph{pre2020} fit estimates on data up to 31 December 2019 and is scored against the
later sample, which contains three held-out episodes fixed in the protocol by the end date of the window: March
to July 2020, February to June 2022 and June to July 2024. Windows within $T$ days of a split edge are dropped.
Neural generators use 2018 to 2019 as a validation block for early stopping or checkpoint selection only;
classical generators are estimated on every day of the fit period. Each (generator, fit) pair is run with five
seeds and produces 2000 windows per seed. Sampling is unconditional: no real window, date or label may be
passed to a generator at sampling time. A generator that simulates one long path cuts it into non-overlapping
windows.

\subsection{Crisis label rule}
\label{sec:label}
The label of a window is a function of its own returns only, and the same function is applied to real and
generated windows. Let $r_{t,i}$ be the return of asset $i$ on day $t$ of the window. Define
\begin{align}
\mathrm{rv} &= \frac{1}{20}\sum_{t=T-19}^{T} \operatorname{sd}_i\,(r_{t,i}), \\
\mathrm{dd} &= \frac{1}{N}\sum_{i=1}^{N} \mathbf{1}\Big[\sum_{t=T-4}^{T} r_{t,i} < -0.05\Big], \\
\mathrm{cp} &= \frac{1}{N}\sum_{i=1}^{N} \mathbf{1}\big[\,|r_{T,i}| > 0.095\,\big],
\end{align}
where $\operatorname{sd}_i$ is the cross-sectional standard deviation. The score is
\begin{equation}
s = 0.5\,\frac{\mathrm{rv}}{\mathrm{rv}_{\mathrm{ref}}} + 0.3\,\mathrm{dd} + 0.2\,\mathrm{cp},
\qquad \text{crisis} \iff s \ge s^{\ast}.
\label{eq:score}
\end{equation}
The reference $\mathrm{rv}_{\mathrm{ref}}$ is the 99th percentile of $\mathrm{rv}$ and the threshold $s^{\ast}$
is the 85th percentile of $s$, both estimated on the real windows of the fit period and then frozen. They are
never re-estimated on generated data, so a generator whose windows are all calm receives a crisis frequency of
zero. The third term was designed as a proximity-to-limit indicator; since the stocks in the panel have no
fixed band it acts only as an indicator of very large single-day moves. The weights and percentiles are
judgment calls made before the runs. A second rule, used for a sensitivity table and never for ranking, labels
a window as crisis when the 20-day realised volatility of the equal-weight portfolio is at or above its 85th
percentile on the fit-period real windows. If a generator produces fewer than 20 crisis windows its crisis
metrics are recorded as not produced.

The rule contains no correlation term. A crisis is defined by dispersion, drawdown breadth and large moves, and
the benchmark then asks what dependence looks like inside the windows so labelled.

\subsection{Metric families}
\label{sec:metrics}
All windows, real, bootstrap or generated, are reduced to per-window statistics by one code path. ``Pooled''
means over all windows and days.

\textbf{Family A, marginal and temporal:} ratio of pooled standard deviations; median excess kurtosis over
assets; lag-1 autocorrelation of returns; lag-1 and mean lag-1-to-5 autocorrelation of squared returns;
skewness of the equal-weight portfolio.

\textbf{Family B, static dependence:} Frobenius norm and mean absolute off-diagonal difference between the
pooled generated and real correlation matrices; $L_2$ distance between their sorted eigenvalue spectra divided
by $\sqrt{N}$; share of the largest eigenvalue.

\textbf{Family C, crisis dependence:} crisis frequency under the frozen rule; the crisis correlation boost
\begin{equation}
\mathrm{boost} = 100\,\frac{\bar m_{\mathrm{crisis}} - \bar m_{\mathrm{normal}}}{\bar m_{\mathrm{normal}}},
\label{eq:boost}
\end{equation}
where $m$ is the mean absolute off-diagonal entry of a window's own 60-day correlation matrix and the bars are
means over windows of each label; tail co-exceedance, the average over ordered pairs of
$P(r_j<q_j \mid r_i<q_i)$ at the pooled 5\% quantiles; the ratio of the 5\% expected shortfall of equal-weight
returns in crisis windows to that in normal windows (last 20 days of each window); and the crisis-minus-normal
gaps of the top-eigenvalue share \cite{kritzman2011absorption} and of the share of pairs with $|\rho|>0.5$
\cite{diebold2014network}.

\textbf{Family D, distinguishability:} accuracy of a gradient-boosted classifier, real against generated, under
blocked cross-validation with a purge of $T$ days, on 20 order-free summary features of a window and on the
first 64 principal components of the flattened window.

\subsection{Noise ceiling and scores}
\label{sec:ceiling}
A finite real sample does not determine its own statistics exactly, and a generator should not be penalised
for a gap smaller than that uncertainty. For each fit we draw 500 circular block bootstrap resamples of the
real window indices with a block of $2T$ windows \cite{politis1994stationary}. For a statistic metric with real
value $\theta$ and 95\% bootstrap interval $[\ell,u]$, the normalised error of a generator with value
$\hat\theta$ is
\begin{equation}
e = \min\!\Big(\frac{|\hat\theta-\theta|}{(u-\ell)/2},\,10\Big).
\label{eq:nerr}
\end{equation}
For a distance metric, $e$ is the generated-to-real distance divided by the 95th percentile of the distance
between a bootstrap resample of the real data and the real data, with the same cap. An error of one therefore
means ``as far from the real value as real data is from itself''. The family score $S_F$ is the mean of the
capped errors of the metrics in family $F$, generators are ranked by $S_F$, and agreement between two rankings
is measured by Kendall's $\tau$ \cite{kendall1938rank}. We report $\tau$ with an interval over seeds and with
its range when each generator is left out in turn.

\subsection{Controlled ground truth}
The boost in \eqref{eq:boost} is a ratio of averages of absolute sample correlations over short windows with
estimated labels. To learn what it measures we build simulated panels with the shape and dates of the real
panel. Four panels follow a two-state Markov-switching one-factor model with Student-t shocks in which the
crisis state multiplies all volatilities and, separately, raises the factor loading so that the population
mean correlation in crisis exceeds that in the normal state by a known amount (0, 20, 40 and 80\%). They share
the state path and the shocks. A fifth panel has no regimes: a constant-conditional-correlation GARCH(1,1) in
which volatility clusters and the conditional correlation matrix never changes, so its true boost is zero. The
unchanged benchmark code is run on these panels.

\subsection{Contamination check}
To test whether metric-based orderings survive data errors (H5), three contaminated copies of the cleaned
panel are scored: one cell replaced by a single bad tick, 1\% of all asset-days replaced by spikes, and one
cell replaced by an unadjusted two-for-one split. For each metric we count the generator pairs whose order
changes relative to the clean panel.

\subsection{Downstream risk validation}
\label{sec:riskdesign}
A crisis metric is useful if the ranking it produces says something about a task. For each out-of-sample day
from January 2020 onward and each generator we forecast the one-day VaR and ES of a portfolio from a pool of
20\,000 generated windows: the state of the day is the realised volatility of the real portfolio over the
previous 20 days, and the forecast is the empirical VaR and ES of the next-day returns of the 5\% of pool
entries whose synthetic 20-day volatility is closest to that state. The pool used in calendar year $Y$ comes
from the generator refitted on data up to the end of year $Y-1$, so every forecast is out of sample. Portfolios
are equal-weight and one fixed random-weight portfolio; levels are 95, 97.5 and 99\%. Benchmarks are rolling
historical simulation over 500 days and GARCH(1,1)-t filtered historical simulation. We report the
unconditional coverage test \cite{kupiec1995techniques}, the independence and conditional coverage tests
\cite{christoffersen1998interval}, the Basel traffic light \cite{bcbs1996backtesting}, the $Z_2$ ES test
\cite{acerbi2014backtesting} and the FZ0 joint loss for VaR and ES
\cite{fissler2016elicitability,patton2019dynamic}. Generators are ranked by mean FZ0 loss at 97.5\% on the
equal-weight portfolio; this choice was frozen. Metric validity is the Kendall $\tau$ between that ranking and
the ranking by each family score from the pre2020 fit, with a joint bootstrap over backtest days (blocks of 20)
and seeds.

\subsection{Preregistered hypotheses and gates}
Table~\ref{tab:prereg} states the hypotheses and gates as frozen. The thresholds are judgment calls made before
the runs and are not evidence in themselves.

\begin{table}[t]
\caption{Hypotheses and gates frozen before any generator was scored.}
\label{tab:prereg}
\centering
\footnotesize
\begin{tabular}{@{}lp{0.86\columnwidth}@{}}
\toprule
H1 & Kendall $\tau$ between the family-B and family-C rankings of generators is below 0.3. \\
H2 & At least one non-neural regime model (regime-conditional Gaussian or HMM) has a family-C score within one
seed standard deviation of CG-NSDE, or better. \\
H3 & Kendall $\tau$ between the family-C rank and the VaR/ES backtest rank exceeds 0.5 and exceeds the $\tau$ of
family B, with non-overlapping bootstrap intervals. \\
H4 & The sign of ($\tau_{B,C}<0.3$) replicates on at least one of the two other markets. \\
H5 & Injecting a single bad tick flips the order of at least one generator pair on at least one metric. \\
\midrule
G1 & Pass if every regime-aware generator beats every static generator on the family-C score and no
regime-conditional Gaussian or HMM matches CG-NSDE (H2 false). Fail: the benchmark leads the paper. \\
G2 & Pass if the held-out boost of CG-NSDE lies inside the real 95\% interval on at least two of three episodes. \\
G3 & Pass if H3 holds. \\
G4 & Pass if H4 holds. \\
Kill & If static generators also rank well on family C across seeds (H1 fails on every market), drop the
benchmark framing. \\
\bottomrule
\end{tabular}
\end{table}

% ===========================================================================
\section{Generators}
\label{sec:generators}
% ===========================================================================
Eleven generators are registered. All are our own implementations on standard libraries.

\textbf{Static.} (i) i.i.d. Gaussian days with the sample mean and covariance. (ii) A static multivariate
Student-t with a rank-8 factor covariance from probabilistic PCA \cite{tipping1999ppca}, one mixing variable
per day and degrees of freedom by profile likelihood. Both have no regime and no volatility clustering by
construction.

\textbf{Resampling.} (iii) The stationary bootstrap \cite{politis1994stationary} of whole cross-sections. It
replays real days and is a strong in-sample reference, not a model that can extrapolate.

\textbf{Dynamic classical.} (iv) GARCH(1,1)-t per asset with filtered historical simulation of whole residual
cross-sections \cite{bollerslev1986garch,baroneadesi1999fhs}. (v) DCC(1,1) with correlation targeting
\cite{engle2002dcc}. Deviations: persistence is clipped below one; DCC parameters are found by Nelder-Mead on a
constrained reparametrisation; DCC simulation uses multivariate Student-t innovations with one common degree
of freedom, which is not the estimated model and adds joint tail dependence.

\textbf{Regime.} (vi) A supervised two-regime Gaussian: each fit-period day takes the frozen crisis label of
the window ending on it, per-regime means and full covariances are sample moments (shrunk to the pooled
covariance when a regime has few days), and days follow the estimated two-state Markov chain
\cite{hamilton1989regime}. Day labels are trailing statistics and lag the regime. (vii) An unsupervised
Gaussian hidden Markov model \cite{rabiner1989hmm} fitted by Baum-Welch on standardised returns, with the
number of states chosen by BIC and a ridge in the M-step.

\textbf{Neural baselines.} (viii) TimeGAN \cite{yoon2019timegan}, (ix) Quant GAN \cite{wiese2020quantgans} with
the Lambert-W preprocessing of \cite{goerg2015lambert}, (x) TimeVAE \cite{desai2021timevae}, base architecture.
These are CPU-scale reimplementations in PyTorch, not the authors' code, and the comparison must be read with
that in mind. TimeGAN is trained for a few hundred iterations per phase where the original uses tens of
thousands; its output is bounded by the per-asset training range because of min-max scaling. Quant GAN is
extended from one series to $N$ channels and is trained with a gradient penalty instead of weight clipping,
with a small fixed iteration budget. TimeVAE uses z-scored inputs and early stopping on the validation
objective, and its decoder mean is sampled without observation noise as in the original, so its windows are
smoother than real ones. Some loss weights of TimeGAN and TimeVAE follow the released code as recalled and were
not re-checked against the source. A weak score of a neural baseline here is evidence about the baseline under
this compute budget and not about the published method.

\textbf{Reference method.} (xi) CG-NSDE, described next.

% ===========================================================================
\section{CG-NSDE: Preregistered Model, Its Failure, and a Revised Model}
\label{sec:method}
% ===========================================================================
Two versions of our own generator appear in the tables. Version 1 is the model that was registered in the
protocol. It failed, for reasons we could diagnose, and version 2 was designed afterwards. Version 2 is
therefore a \emph{post hoc} model: its design used knowledge of the version 1 results, and every statement
about it is exploratory in the sense of the protocol. Both are scored by the unchanged benchmark and both stay
in every table. All models work on per-asset z-scored returns $x_{1:T}\in\mathbb{R}^{T\times N}$, with the
scaler estimated on training days only.

\subsection{Shared components}
\textbf{Graph encoder.} For each training window a symmetric $k$-nearest-neighbour graph ($k=5$, with self
loops) is built on the absolute Pearson correlation of the window. Each node carries six features: mean, mean
absolute value and standard deviation of its z-scored returns, mean z-scored log volume, a sector code and its
largest absolute raw return in the window. A two-layer graph attention network \cite{velickovic2018gat} (four
heads, then one) followed by mean pooling gives a vector $g$.

\textbf{Latent SDE.} A GRU encoder gives a diagonal Gaussian posterior $q(z_0\mid x)$; the prior
$p(z_0\mid y)$ is a learned diagonal Gaussian for each regime $y\in\{0,1\}$. The state evolves on $[0,1]$ as
\begin{equation}
dz_t = \tanh\!\big(f_\theta(z_t,c,t)\big)\,dt
     + \sigma_{\max}\,\operatorname{sigm}\!\big(h_\theta(z_t,c,t)\big)\odot dW_t ,
\label{eq:sde}
\end{equation}
with multilayer perceptrons $f_\theta$, $h_\theta$, a context $c$ and diagonal noise, solved by the Euler
scheme with $T$ fixed steps; day $t$ is emitted from the state after $t$ steps. This is not the path-space
latent SDE of \cite{li2020scalable}: only $z_0$ has a posterior, and the Brownian path is drawn from the prior
dynamics during training \cite{kidger2021neuralsde}.

\textbf{Auxiliary losses.} $\mathcal{L}_{\mathrm{con}}$ is the supervised contrastive loss
\cite{khosla2020supcon} on a projection of the time-averaged latent path, with the regime as the class.
$\mathcal{L}_{\mathrm{graph}}$ is the squared distance between each node's second-layer attention row and the
row-normalised absolute correlation over its neighbours. The training labels $y$ come from the frozen rule
applied to the training windows. They enter the prior, the contrastive loss and the batch sampler (batches of
32 windows hold exactly six crisis windows), and they are not available at sampling time.

\textbf{Unconditional sampling.} After training, every training window is encoded and a diagonal Gaussian
mixture is fitted to the latent codes of each regime, with the number of components chosen by BIC
\cite{ghosh2020rae}. To sample, a regime is drawn with the training crisis frequency, a code is drawn from
that regime's mixture, the SDE is solved and returns are emitted day by day. No real window, graph or label is
used. A generated window is labelled crisis only if the frozen rule of Section~\ref{sec:label} says so,
whatever regime was drawn.

\subsection{Version 1 (preregistered)}
In version 1 the context of the SDE is the graph vector, $c=g\in\mathbb{R}^{64}$, and $z_0\in\mathbb{R}^{64}$.
Linear heads on $z_t$ give a mean $m_t$, a diagonal $d_t>0$ and loadings $B_t\in\mathbb{R}^{N\times K}$, $K=8$,
and the day is multivariate Student-t with one mixing variable,
\begin{equation}
x_t \mid z_t, v_t \sim t_\nu\!\big(m_t,\; \lambda_t\,(B_tB_t^{\top} + \operatorname{diag} d_t)\big),
\label{eq:emis}
\end{equation}
where $\lambda_t=\exp(2\kappa v_t)$ and the scalar state $v_t$ is an exponentially weighted average of the
previous days' mean squared standardised shocks. The loss is the negative log likelihood plus
$\beta\,\mathrm{KL}(q\|p)$ with $\beta=0.01$ (KL averaged over dimensions), a pooled-correlation matching term
with weight 100, and the two auxiliary losses with weight 0.1. The mixture is fitted to $[z_0,g]$.

\textbf{Why it failed.} Two design errors are visible in the run logs and in the equations. (i) The decoder
receives 128 numbers per window, $z_0$ and $g$, that are computed from the window it has to reconstruct and are
almost unregularised. With training windows that overlap in 59 of 60 days the model memorises: the training
likelihood improves for as long as training runs, while the validation likelihood is best after
\vOneBestEpoch\ of \vOneEpochs\ epochs, so early stopping selects a model whose dynamics parameters have barely
moved from their initial values. (ii) The volatility state multiplies the whole covariance matrix by one
scalar. This changes variances and leaves every correlation unchanged, so the only mechanism that reacts to
past returns cannot produce a rise of correlation under stress. The generated windows are close to
independent multivariate Student-t days.

\subsection{Version 2 (post hoc): factor-volatility states}
Version 2 keeps the graph encoder, the regime prior, the latent SDE and both auxiliary losses, and changes
where information flows and what the volatility state can do.

\textbf{Information bottleneck.} The graph vector enters only the posterior,
$q(z_0\mid x)=\mathcal{N}\big(\mu(\mathrm{GRU}(x),g),\operatorname{diag}\sigma^2(\cdot)\big)$ with
$z_0\in\mathbb{R}^{16}$, and the SDE has no context. Everything the decoder knows about a window passes through
$z_0$, whose KL divergence to the regime prior is summed over dimensions with weight one. The objective is the
per-day negative evidence lower bound plus the auxiliary terms,
\begin{equation}
\mathcal{L} = \mathcal{L}_{\mathrm{NLL}} + \tfrac{1}{T}\,\mathrm{KL}\big(q(z_0\mid x)\,\|\,p(z_0\mid y)\big)
 + \lambda_{s}\mathcal{L}_{\mathrm{con}} + \lambda_{g}\mathcal{L}_{\mathrm{graph}},
\label{eq:loss}
\end{equation}
and early stopping uses the same bound on the validation block. The correlation-matching term is dropped.

\textbf{Emission.} Loadings $B\in\mathbb{R}^{N\times K}$ and idiosyncratic variances $d\in\mathbb{R}^N_{>0}$
are free parameters, initialised from the leading $K=8$ principal components of the training correlation
matrix. The covariance of day $t$ is
\begin{equation}
\Sigma_t = B\,\operatorname{diag}\!\big(e^{\,v^{f}_t + a^{f}(z_t)}\big)B^{\top}
         + \operatorname{diag}\!\big(d\odot e^{\,v^{i}_t + a^{i}(z_t)}\big),
\label{eq:sigma}
\end{equation}
with $x_t\mid z_t,v_t\sim t_\nu(0,\Sigma_t)$. Here $v^{f}_t\in\mathbb{R}^K$ are log-variance states of the
factors, $v^{i}_t\in\mathbb{R}^N$ are log-variance states of the idiosyncratic terms, and
$a(z_t)=1.5\tanh(Wz_t)\in\mathbb{R}^{K+1}$ is a shift read from the latent path (one value per factor and one
common idiosyncratic value). When the factor variances rise relative to the idiosyncratic variances, pairwise
correlations rise. Correlation change is therefore produced by the dynamics and does not have to be stored in
a latent code.

\textbf{State recursion.} Let $\hat f_t=\mathbb{E}[f_t\mid x_t]$ be the posterior mean of the standardised
factors under $\Sigma_t$, $c_{t,k}=\operatorname{Var}(\hat f_{t,k})$ and $\hat\varepsilon_t$ the
idiosyncratic residual with variance $s_{t,i}$; all three follow in closed form from the $K\times K$ matrix
already used for the likelihood. Then
\begin{align}
v^{f}_{t+1,k} &= \rho_k v^{f}_{t,k} + \alpha_k\Big(\frac{\hat f_{t,k}^2}{c_{t,k}}-1\Big)
                + \theta_k\frac{\hat f_{t,k}}{\sqrt{c_{t,k}}}, \label{eq:vf}\\
v^{i}_{t+1,i} &= \rho_0 v^{i}_{t,i} + \alpha_0\Big(\frac{\hat\varepsilon_{t,i}^2}{s_{t,i}}-1\Big), \label{eq:vi}
\end{align}
with shocks capped at 20 and states clipped to $[-3,3.5]$, and $v_1=2\tanh(Vz_0)$. This is an exponential
GARCH recursion \cite{nelson1991egarch} on shocks that are standardised by the current model variance: under
the model the shock terms have mean zero, so the states are stable for any parameter value. $\theta_k$ is a
leverage term. The states are teacher-forced on real returns in training and fed with the model's own draws in
sampling. The scalar parameters $\rho,\alpha,\theta,\nu$ sit behind maps with small slopes and are trained
with a learning rate of $2\times10^{-2}$; all other parameters use $10^{-3}$.

\textbf{Innovations.} By default version 2 samples by filtered historical simulation
\cite{baroneadesi1999fhs}: the standardised residual $\eta=\Sigma_t^{-1/2}x_t$ (symmetric root) of the last
day of every training window is stored, whole residual cross-sections are resampled, and a generated day is
$\Sigma_t^{1/2}\eta$. The parametric alternative draws from the fitted Student-t and is reported as an
ablation. The residual pool is the only training data kept by the fitted model.

\textbf{Development protocol.} All design choices of version 2 were made on the 2010 to 2019 fit period: the
model was fitted on its training windows and its samples were compared with the real windows of the same
period. Four configurations were tried in sequence (network-valued loadings with the network learning rate;
the same with a separate rate for the dynamics parameters; static loadings with per-asset states; the same
with resampled innovations) and the last was frozen before any benchmark job of version 2 was started. The
2020 to 2024 held-out sample, the other markets and the risk backtest were not consulted. The log of these
runs, including the two that failed, is in the repository.

\textbf{Ablations.} Variants of version 2 remove one component each: the graph attention network (per-node
MLP), the contrastive loss, the graph loss, the SDE (constant latent state), the leverage term, the factor
structure of the states (one scalar state as in version 1), the state recursion altogether, the Student-t
likelihood (Gaussian), the resampled innovations (parametric draws); further variants set the KL weight to
0.1 and the number of factors to 1, 4 and 16.

% ===========================================================================
\section{Experimental Protocol}
% ===========================================================================
The protocol file was frozen on 12 November 2026 (version 1.0.0) and is dumped to the results directory before
scoring; any later change must be logged. All runs use one laptop CPU without a GPU. A resumable job queue runs
every (market, generator, fit, seed) job; failed jobs are logged with their error and are listed with the
results. Tables, figures and the numbers quoted in the text are produced by one build script from the run
outputs.

Deviations from the frozen protocol that we are aware of are the following. The ablation variants are run on
the in-sample fit with three seeds instead of five, for compute reasons. The risk backtest uses the pool of one
seed per refit. The controlled-ground-truth validation uses 200 bootstrap resamples instead of 500.
% RESULT-DEPENDENT: list here any (market, generator, fit) cells that are missing from the final tables because
% jobs failed (at the time of writing: Korean GARCH-FHS pre2020 jobs), with the reason from the queue log.

% ===========================================================================
\section{Results}
\label{sec:results}
% ===========================================================================
\emph{Status of this draft: the tables and figures below are generated from the runs completed so far. Rows for
generators whose jobs have not finished are absent, and all verdicts are provisional until every row is
present.}

\subsection{What the real data supports}
On the main panel the crisis correlation boost is \realNseBoost\% with a 95\% block-bootstrap interval of
[\realNseBoostLo, \realNseBoostHi]. The corresponding values are \realUsBoost\% [\realUsBoostLo,
\realUsBoostHi] for the United States and \realKrBoost\% [\realKrBoostLo, \realKrBoostHi] for Korea. These
intervals are wide: fourteen years of daily data contain few independent crisis episodes, and overlapping
windows do not add independent information. This limits how finely any benchmark can separate generators on
crisis statistics, and it is the reason every error is expressed relative to the interval half-width in
\eqref{eq:nerr}. A generator inside the band cannot be said to be better than another generator inside the band.
% NEED-MACRO: real held-out (2020-24) boost with its interval, and the per-episode real boosts with intervals
% (covid_2020, shock_2022, mid_2024); they exist in results/tables/T3_heldout.csv but not in numbers.tex.

\subsection{Does the boost statistic measure what it claims?}
\label{sec:dgp}
Fig.~\ref{fig:dgp} shows the benchmark's boost on the simulated panels against the known population boost.
Three facts follow, and they constrain every later statement. First, the statistic is monotone in the truth
but strongly compressed: the largest simulated boost is measured well below its population value. Part of
the compression remains even with perfect labels and windows that lie entirely in one state, because the mean
absolute sample correlation of a 60-day window has a positive floor that weighs more when the true correlation
is lower; the rest comes from windows that mix states. Second, when volatility rises in the crisis state but
correlation does not change, the measured boost is close to zero, with an interval that covers zero. Third,
and less comfortably, on the GARCH panel without any regime the measured boost is clearly positive and its
interval excludes zero. Volatility clustering alone therefore produces a positive reading, in line with the
heteroskedasticity argument of \cite{forbes2002contagion}. The boost in \eqref{eq:boost} is an ordinal
statistic, biased towards zero as a measure of the population boost and biased away from zero under pure
volatility clustering. It can be compared between a generator and real data under the same rule; it cannot be
read as a population correlation increase, and a positive value is not proof of a correlation regime. The
validation uses one realisation per panel and a one-factor model in which labelling is easier than on real
data.
% NEED-MACRO: measured boost and interval on sim_boost80 (true 80) and on sim_ccc (no-regime GARCH null, about
% 14 percent per results/01_dgp/README.md), and the 60-day floor of mean |sample correlation| for independent assets.

\begin{figure}[t]
\centering
\maybefig{\columnwidth}{figures/f5_dgp.pdf}
\caption{Controlled ground truth: measured window boost against the known population boost on simulated panels,
with oracle labels and with the frozen rule, and the reading on a GARCH panel without regimes.}
\label{fig:dgp}
\end{figure}

\subsection{In-sample ranking: correlation fidelity against crisis dependence}
Table~\ref{tab:main} reports, for the full fit on the main panel, selected metrics and the four family scores.
Bold marks a value inside the real 95\% band; underlining marks the best family score. Fig.~\ref{fig:scatter}
plots the correlation error against the boost, and Fig.~\ref{fig:scores} shows the family-B and family-C scores
with seed standard deviations.

The two static generators have a boost of \boostIidGaussian\% and \boostTPpca\%, as their construction implies:
with i.i.d. days the only source of a non-zero value is the sampling noise of the labelled subsets.
% RESULT-DEPENDENT: describe which generators lie inside the real boost band and inside the correlation ceiling;
% state whether any generator is inside both. Do not rank CG-NSDE before its 40-epoch runs are in the table.
% RESULT-DEPENDENT: comment on the stationary bootstrap. It replays real days, so a good in-sample family-C score
% is expected by construction; say whether it keeps that position in the held-out fit.
% RESULT-DEPENDENT: comment on GARCH-FHS and DCC relative to the regime models on family C. In the partial
% results the GARCH models, which have no regime, score at least as well on family C as the HMM; together with
% Section VII-B this would mean that part of family C is driven by volatility clustering. Confirm or revise.

The Kendall $\tau$ between the family-B and the family-C ranking over the \nGenerators\ generators scored so far
is \tauBCnseFull\ in sample and \tauBCnseHeld\ in the held-out fit; Table~\ref{tab:taus} gives the intervals
over seeds and the leave-one-generator-out range for each market and fit.
% RESULT-DEPENDENT: H1 verdict sentence (currently generated as \verdictH1 on 7 classical generators only).
% State the final tau, its interval, the leave-one-out range, and whether the interval excludes 0.3. With about
% eleven generators a Kendall tau has a wide sampling distribution: say so whatever the verdict.

\begin{table*}[t]
\caption{Main panel (India), full fit: selected metrics and family scores $S_A$ to $S_D$ (mean normalised error;
one means as far from the real value as real data is from itself). Bold: inside the real 95\% band.
Underlined: best family score.}
\label{tab:main}
\centering
\footnotesize
\resizebox{\textwidth}{!}{\maybetable{tables/t1_nse_full.tex}}
\end{table*}

\begin{figure}[t]
\centering
\maybefig{\columnwidth}{figures/f1_scatter_nse_full.pdf}
\caption{Correlation error against crisis correlation boost, full fit. Horizontal band: real 95\% interval of
the boost. Vertical band: real-versus-real ceiling of the correlation error.}
\label{fig:scatter}
\end{figure}

\begin{figure*}[t]
\centering
\maybefig{0.9\textwidth}{figures/f2_scores_nse_full.pdf}
\caption{Family-B and family-C scores, full fit, with seed standard deviations. The shaded band is an error of
at most one unit of real sampling noise.}
\label{fig:scores}
\end{figure*}

\begin{table}[t]
\caption{Kendall $\tau$ between the family-B and family-C rankings.}
\label{tab:taus}
\centering
\footnotesize
\resizebox{\columnwidth}{!}{\maybetable{tables/t6_taus.tex}}
\end{table}

\subsection{Held-out crises}
Table~\ref{tab:heldout} gives the same quantities for the pre2020 fit, scored on 2020 to 2024, and
Fig.~\ref{fig:heldout} shows the boost of each generator against the real value and interval for each held-out
episode. The real intervals per episode are wider than the in-sample one because each episode contains few
crisis windows. The held-out fit is a harder and different test from the in-sample one: a generator estimated
up to 2019 has not seen a shock of the size of March 2020.
% RESULT-DEPENDENT: say how many generators fall inside the real interval for each episode. In the partial
% results every classical generator is far below the real COVID-2020 interval; confirm for the final set and
% state it plainly, including for CG-NSDE (gate G2).
% RESULT-DEPENDENT: say whether the family-B versus family-C disagreement is stronger, weaker or of different
% sign in the held-out fit than in sample (\tauBCnseHeld versus \tauBCnseFull).

\begin{table*}[t]
\caption{Main panel (India), pre2020 fit scored on the held-out 2020 to 2024 sample. Layout as in
Table~\ref{tab:main}.}
\label{tab:heldout}
\centering
\footnotesize
\resizebox{\textwidth}{!}{\maybetable{tables/t1_nse_pre2020.tex}}
\end{table*}

\begin{figure*}[t]
\centering
\maybefig{0.95\textwidth}{figures/f3_heldout.pdf}
\caption{Held-out episodes: crisis correlation boost of each generator (pre2020 fit, mean and seed standard
deviation) against the real value (line) and its 95\% interval (band). $n$ is the number of real crisis
windows in the episode.}
\label{fig:heldout}
\end{figure*}

\subsection{Downstream validation: VaR and ES backtest}
Table~\ref{tab:risk} reports the backtest of the equal-weight portfolio at the 99\% level for each generator
and for the two benchmarks, sorted by FZ0 loss, and Fig.~\ref{fig:var} shows the forecast paths.
Table~\ref{tab:validity} gives the Kendall $\tau$ between the FZ0 ranking and the ranking by each family score.
% RESULT-DEPENDENT: describe coverage: which methods are in the green, yellow or red zone, and whether any
% generator-based forecast is competitive with historical simulation and GARCH-FHS.
% RESULT-DEPENDENT: H3 and gate G3. State tau for family C (\tauRiskC) and family B (\tauRiskB) with intervals,
% and whether the intervals overlap. If H3 fails, say that the crisis family is not validated by this task.
The backtest covers one out-of-sample period of five years on one market with one seed per refit, and the
nearest-neighbour conditioning on a single volatility state is a simple choice; a different conditioning could
change the ranking.

\begin{table}[t]
\caption{VaR/ES backtest, equal-weight portfolio, 99\% level, out of sample from 2020.}
\label{tab:risk}
\centering
\footnotesize
\resizebox{\columnwidth}{!}{\maybetable{tables/t4_risk.tex}}
\end{table}

\begin{table}[t]
\caption{Metric validity: Kendall $\tau$ between the backtest ranking (FZ0 at 97.5\%) and the ranking by each
benchmark score.}
\label{tab:validity}
\centering
\footnotesize
\maybetable{tables/t5_validity.tex}
\end{table}

\begin{figure*}[t]
\centering
\maybefig{0.95\textwidth}{figures/f4_var.pdf}
\caption{Daily returns of the equal-weight portfolio and one-day 99\% VaR forecasts, out of sample.}
\label{fig:var}
\end{figure*}

\subsection{Replication on other markets}
Table~\ref{tab:markets} repeats the in-sample table for the United States and Korea. The in-sample
$\tau$ between the family-B and family-C rankings is \tauBCusFull\ for the United States and \tauBCkrFull\ for
Korea (Table~\ref{tab:taus}).
% RESULT-DEPENDENT: H4 and gate G4. Say on which market or markets tau is below 0.3 and on which it is not. In
% the partial results the US value is above the 0.3 threshold and the Korean value is below it; H4 needs only
% one market, so report both values and do not describe this as replication on both.
CG-NSDE uses a sector feature defined for the Indian tickers only; on the other markets that feature is zero.
% NEED-MACRO: an out-of-time check on the Indian panel extended to October 2026 exists as
% tables/t1_nse_ext_upto2024.tex once those jobs finish; add a sentence and the table if space allows.

\begin{table*}[t]
\caption{Replication panels, full fit. Top: United States. Bottom: Korea. Layout as in Table~\ref{tab:main}.}
\label{tab:markets}
\centering
\footnotesize
\resizebox{\textwidth}{!}{\maybetable{tables/t1_us_full.tex}}\\[4pt]
\resizebox{\textwidth}{!}{\maybetable{tables/t1_kr_full.tex}}
\end{table*}

\subsection{CG-NSDE ablations}
Table~\ref{tab:ablation} and Fig.~\ref{fig:ablation} compare the full model with its variants on the in-sample
fit (three seeds).
% RESULT-DEPENDENT: for each removed component say whether the family-C and family-A scores move by more than
% one seed standard deviation. A component whose removal changes nothing must be reported as not supported by
% the ablation. Do not claim that any component is necessary unless the table shows it.
% RESULT-DEPENDENT: H2 and gate G1. Compare CG-NSDE with the regime-conditional Gaussian and the HMM on family C
% (\boostCgnsde and the S_C column). If a classical regime model matches it, say so.

\begin{table}[t]
\caption{CG-NSDE ablations, main panel, full fit.}
\label{tab:ablation}
\centering
\footnotesize
\resizebox{\columnwidth}{!}{\maybetable{tables/t2_ablation.tex}}
\end{table}

\begin{figure*}[t]
\centering
\maybefig{0.9\textwidth}{figures/f6_ablation.pdf}
\caption{Ablations: family-C and family-A scores of the full model and its variants, with seed standard
deviations.}
\label{fig:ablation}
\end{figure*}

\subsection{Sensitivity to data errors}
The contamination check counts, for each metric, the generator pairs whose order changes when the benchmark is
run on a contaminated copy of the panel.
% NEED-MACRO: number of metrics with at least one pair flip, and the largest flip count, for each of the three
% contaminated panels (results/tables/T7_contamination.csv has no LaTeX table and no macros yet).
% RESULT-DEPENDENT: H5 verdict (currently \verdictH5 on the classical generators). Note that many flips under a
% single bad tick involve pairs whose scores were already within seed noise of each other; say whether any flip
% involves a pair separated by more than one seed standard deviation.

\subsection{Hypotheses and gates}
Table~\ref{tab:verdicts} lists the verdict that the analysis code assigns to each hypothesis and gate, with the
evidence string it generates. Items without a row have not yet been evaluated.
% RESULT-DEPENDENT: one paragraph summarising which of H1-H5 and G1-G4 passed and failed, stating failed items
% first, and what the paper can and cannot claim as a consequence (per the gate texts in Table I).

\begin{table*}[t]
\caption{Verdicts on the preregistered hypotheses and gates, as generated by the analysis code.}
\label{tab:verdicts}
\centering
\footnotesize
\maybetable{tables/t8_hypotheses.tex}
\end{table*}

% ===========================================================================
\section{Threats to Validity}
\label{sec:threats}
% ===========================================================================
\textbf{Few crises.} The sample holds a small number of independent stress episodes. The real intervals are
wide, the held-out episodes are three, and many generator differences on family C lie inside the noise.

\textbf{The boost statistic.} As Section~\ref{sec:dgp} shows, the boost is compressed and gives a positive
reading under volatility clustering without any correlation regime. A generator with good volatility dynamics
and constant conditional correlation can therefore obtain part of the real boost. Family C measures
``dependence in windows that look stressed'', which is what a stress tester observes, but it does not identify
the mechanism.

\textbf{Label rule.} The rule is one of many reasonable ones, and its weights, percentile and lookbacks are
judgment calls. It labels by a relative threshold, so by construction about 15\% of real fit-period windows are
crisis. The supervised regime Gaussian and CG-NSDE are trained with labels from the same rule that later
scores them, which may favour them on crisis frequency; the unsupervised models have no such advantage. One
term of the rule and one node feature of CG-NSDE were designed around a 10\% price limit that does not apply
as a fixed band to the stocks in the panel.

\textbf{Data.} Survivorship bias removes failed and delisted firms and probably understates tail events. Yahoo
data may contain adjustment errors that the bad-tick rule does not catch. The contamination check shows how
sensitive pairwise orderings are to such errors.

\textbf{Baselines.} The neural baselines are reimplementations trained on a CPU with budgets far below the
original papers, with stated deviations. The classical baselines also deviate from textbook versions in the
ways listed in Section~\ref{sec:generators}. CG-NSDE received more of our development time than any baseline.

\textbf{Statistics.} Kendall's $\tau$ over about eleven generators is coarse. The thresholds in the hypotheses
(0.3 and 0.5) were chosen before the runs but are arbitrary. Family scores average capped errors of correlated
metrics with equal weight. Seeds capture training and sampling noise, not data uncertainty, which enters only
through the real-data bootstrap.

\textbf{Scope.} Daily returns of about fifty large stocks per market, windows of 60 days. Nothing here speaks
to intraday data, to other asset classes or to much larger universes.

\textbf{Novelty claim.} The related-work comparison is based on abstracts and a limited search; a full-text
review may find prior work that contains one or more of the elements we combine.

% ===========================================================================
\section{Reproducibility}
% ===========================================================================
The protocol, the panel builder with its audit log and hashes, the single metric module, all generators, the
job queue, the analysis code and the script that builds every table, figure and number in this paper are in one
repository. Raw vendor prices are not redistributed; the panel builder downloads them again and records the
hash of the cleaned matrix so that differences from our panel can be detected. Each run stores its metrics, its
seed and its training information. Unit tests cover the metric pipeline, the generators and the risk
statistics.

% ===========================================================================
\section{Ethics and Intended Use}
% ===========================================================================
The data are public market prices and contain no personal information. Synthetic scenarios produced by any of
the generators studied here are research output and are not validated for regulatory capital, trading or
investment decisions. In particular, the held-out results show how far generated crisis dependence can be from
what occurred, and a stress test built on a generator that passes correlation-fidelity checks alone could
understate joint losses. Nothing in this paper is investment advice.

% ===========================================================================
\section{Conclusion}
% ===========================================================================
We set out to test, under a protocol frozen in advance, whether fidelity to the unconditional correlation
matrix and fidelity to crisis-time dependence order synthetic return generators in the same way, and whether
either ordering is informative for a VaR/ES task. The benchmark built for this test expresses every error in
units of real sampling noise, labels generated windows by their own returns, and checks its main statistic on
data with known ground truth. That check showed limits of the statistic that we report as part of the result:
the crisis correlation boost is ordinal, compressed and sensitive to volatility clustering, and the real-data
intervals are wide.
% RESULT-DEPENDENT: two or three sentences stating the final verdicts on H1, H3 and H4, the position of CG-NSDE
% and of the classical regime models (H2, G1, G2), and which claims are withdrawn because a gate failed.
Future work should use point-in-time constituent data, longer histories with more crises, a crisis statistic
that separates correlation change from volatility change, and baselines trained at their published scale.

\bibliographystyle{IEEEtran}
\bibliography{references}

\end{document}
